\section{Discussion}
\label{sec:discussion}

\paragraph{Limitations.}
Two chromosomes of one cohort, and current foundation models are larger than the 117M we audit,
where capacity is known to matter \cite{morris2025capacity}; our own scale control reaches 59.9M.
Equal supervised tokens is one definition of matched compute and equal attention cost is another,
and we report the first with the wall clock beside it. The character-level result is a statement about matched rounds and about the utility optimum, not an
impossibility claim. On chromosome~22 the released $k$-mer family is not ordered by $k$ under either
normalisation and the two disagree, so we base no ordering on that dataset
(Sec.~\ref{sec:published}); its thresholded extraction rate is not monotone in $k$ either, and we
have no mechanism for that. Our canaries are contiguous haplotype segments, and whether the rarity
mechanism predicts leakage for records of a different shape, such as sparse variant panels, is open.
Finally, the size of the large-vocabulary excess is sensitive to how the embedding is parametrised,
for an identical function class; the character-to-\kmer{3} step and the objective effect are not.

\paragraph{Future work.}
Parameter-efficient fine-tuning interacts with our mechanism in a way we cannot predict: LoRA
restricts which parameters can absorb a record, and the tokeniser determines how rare the tokens are
that would have to absorb it. The mechanism also suggests a predictor that needs no training --- the
corpus frequency of the tokens spelling a record is computable from the tokeniser and the corpus
alone --- which would let an auditor rank candidate tokenisers before spending a GPU hour. And
differential privacy and these settings are complementary rather than alternative; the right
comparison is at matched utility, which needs its own study.

\paragraph{Ethics and data.}
Every input is public. 1000~Genomes phase~3 is consented for open distribution, and the canaries are
built from released genotypes rather than any private cohort. We report aggregate recovery rates and
publish neither reconstructed haplotypes nor carrier identifiers. The attack needs only query access
to a fine-tuned model, which is what a released checkpoint grants, and we describe it so that groups
releasing checkpoints trained on cohort data can test for it before release.

\paragraph{Conclusion.}
Memorisation in a genomic language model is not an accident of a particular release: it follows from
the training objective and from how finely DNA is cut, both chosen before training begins. It is
invisible to every utility measurement a practitioner would normally take --- the held-out-loss
optimum sits one snapshot from complete reproduction, and six standard downstream tasks do not move
across that boundary --- and on real human sequence what comes back is the carrier's rare alleles,
for records that appear once. The two settings that reduce it are available in advance and, in our
experiments, cost nothing in held-out accuracy: train with a masked objective wherever generative
output is not required, and tokenise as finely as the compute budget allows.
